% --- Archon physics formalization source begin ---
% archon:physics
% archon:covers PhyXMiniProblems/problem_phyx_mini_0738.lean
% archon:source-report reports/phyx_mini/problem_phyx_mini_0738.source.json

\chapter{Physics problem phyx\_mini\_0738}
\label{ch:PhyXMiniProblems_problem_phyx_mini_0738}

\paragraph{Problem source.}
\#\# Physical scenario

A certain airplane has a speed of \$290.0\textbackslash{} km/h\$ and is diving at an angle of \$\textbackslash{}theta = 30.0\textasciicircum{}\textbackslash{}circ\$ below the horizontal when the pilot releases a radar decoy. The horizontal distance between the release point and the point where the decoy strikes the ground is \$d = 700\textbackslash{} m\$.

\#\# Question

How high was the release point?

\#\# Answer choices

- A: A: 619 m

- B: B: 751 m

- C: C: 897 m

- D: D: 926 m

\#\# Figure caption (auxiliary; use the image as primary evidence)

The image depicts a diagram featuring an aircraft in a nose-down dive at an angle \textbackslash{}(\textbackslash{}theta\textbackslash{}). There's a trajectory line showing the path of the aircraft, which curves towards a horizontal surface that seems to represent the ground. The aircraft is positioned along this trajectory.

There are two arrows along the path indicating the direction of motion, starting from above and moving down towards and then horizontally along the surface. Additionally, there is a horizontal line from the point of impact extending toward the left labeled with \textbackslash{}(d\textbackslash{}), indicating a distance from the impact point to some reference point. 

The ground is represented as a textured horizontal line at the bottom of the image.

\#\# Dataset metadata

Kinematics; Physical Model Grounding Reasoning, Spatial Relation Reasoning

\paragraph{Recorded answer/context.}
C: C: 897 m

\paragraph{Figure/image path.}
/root/proposal\_for\_physic/science-mango/phyx\_mini\_run/phyx\_data/test\_image/738.png

\paragraph{Formalization target.}
create a compiling Lean file with sorry bodies at `PhyXMiniProblems/problem\_phyx\_mini\_0738.lean`. The Lean declarations must preserve the physical quantities, dimensions or dimensional roles, figure labels, governing-law hypotheses, and final relation expressed by this problem.
Use Mathlib/Physlib names found through LeanExplore where available. If a physics API is missing, introduce faithful local abstractions rather than scalar placeholder aliases.

\begin{theorem}[Physics formalization target]
\label{thm:physics:phyx_mini_0738:target}
\lean{PhyXMiniProblems.ProblemPhyXMini0738.problem_phyx_mini_0738}
\uses{def:physics:phyx-mini-0738:phyxminiproblems-problemphyxmini0738-lengthinmeters, def:physics:phyx-mini-0738:phyxminiproblems-problemphyxmini0738-radardecoyreleasesetup, def:physics:phyx-mini-0738:phyxminiproblems-problemphyxmini0738-matchesprimaryairplanedecoyfigure, def:physics:phyx-mini-0738:phyxminiproblems-problemphyxmini0738-matchesproblemdescription, def:physics:phyx-mini-0738:phyxminiproblems-problemphyxmini0738-usesstandardnearearthgravity, def:physics:phyx-mini-0738:phyxminiproblems-problemphyxmini0738-hasphysicalradardecoyparameters, def:physics:phyx-mini-0738:phyxminiproblems-problemphyxmini0738-satisfiespassivereleaseprojectilelaws, def:physics:phyx-mini-0738:phyxminiproblems-problemphyxmini0738-computedreleaseheightinmeters, def:physics:phyx-mini-0738:phyxminiproblems-problemphyxmini0738-recordeddatasetanswer, def:physics:phyx-mini-0738:phyxminiproblems-problemphyxmini0738-matchesdisplayedwholemeterheight, def:physics:phyx-mini-0738:phyxminiproblems-problemphyxmini0738-isuniqueclosestdisplayedheight}
The assigned autoformalize agent should translate this physics problem into Lean declarations in the covered file, with theorem and lemma proof bodies written as `by sorry`.
\end{theorem}
\begin{proof}
This is an autoformalization task, not a proof task. Produce statements that can later be proved by the physics prover without weakening the physical model.
\end{proof}

% --- Archon declaration topology begin ---
\section{Declaration topology}
\begin{definition}[velocity Dimension]
  \label{def:physics:phyx-mini-0738:phyxminiproblems-problemphyxmini0738-velocitydimension}
  \lean{PhyXMiniProblems.ProblemPhyXMini0738.velocityDimension}
  The physical dimension of velocity, L T⁻¹
\end{definition}
\begin{proof}
  This is the declared model definition of velocity Dimension.
\end{proof}
\begin{definition}[acceleration Dimension]
  \label{def:physics:phyx-mini-0738:phyxminiproblems-problemphyxmini0738-accelerationdimension}
  \lean{PhyXMiniProblems.ProblemPhyXMini0738.accelerationDimension}
  The physical dimension of acceleration, L T⁻²
\end{definition}
\begin{proof}
  This is the declared model definition of acceleration Dimension.
\end{proof}
\begin{definition}[Length Quantity]
  \label{def:physics:phyx-mini-0738:phyxminiproblems-problemphyxmini0738-lengthquantity}
  \lean{PhyXMiniProblems.ProblemPhyXMini0738.LengthQuantity}
  A nonnegative physical length
\end{definition}
\begin{proof}
  This is the declared model definition of Length Quantity.
\end{proof}
\begin{definition}[Signed Length Quantity]
  \label{def:physics:phyx-mini-0738:phyxminiproblems-problemphyxmini0738-signedlengthquantity}
  \lean{PhyXMiniProblems.ProblemPhyXMini0738.SignedLengthQuantity}
  A signed Cartesian coordinate or displacement
\end{definition}
\begin{proof}
  This is the declared model definition of Signed Length Quantity.
\end{proof}
\begin{definition}[Duration Quantity]
  \label{def:physics:phyx-mini-0738:phyxminiproblems-problemphyxmini0738-durationquantity}
  \lean{PhyXMiniProblems.ProblemPhyXMini0738.DurationQuantity}
  A nonnegative physical duration elapsed since release
\end{definition}
\begin{proof}
  This is the declared model definition of Duration Quantity.
\end{proof}
\begin{definition}[Signed Velocity Quantity]
  \label{def:physics:phyx-mini-0738:phyxminiproblems-problemphyxmini0738-signedvelocityquantity}
  \lean{PhyXMiniProblems.ProblemPhyXMini0738.SignedVelocityQuantity}
  \uses{def:physics:phyx-mini-0738:phyxminiproblems-problemphyxmini0738-velocitydimension}
  A signed component of planar velocity
\end{definition}
\begin{proof}
  This is the declared model definition of Signed Velocity Quantity.
\end{proof}
\begin{definition}[Acceleration Quantity]
  \label{def:physics:phyx-mini-0738:phyxminiproblems-problemphyxmini0738-accelerationquantity}
  \lean{PhyXMiniProblems.ProblemPhyXMini0738.AccelerationQuantity}
  \uses{def:physics:phyx-mini-0738:phyxminiproblems-problemphyxmini0738-accelerationdimension}
  A nonnegative magnitude of gravitational acceleration
\end{definition}
\begin{proof}
  This is the declared model definition of Acceleration Quantity.
\end{proof}
\begin{definition}[length In Meters]
  \label{def:physics:phyx-mini-0738:phyxminiproblems-problemphyxmini0738-lengthinmeters}
  \lean{PhyXMiniProblems.ProblemPhyXMini0738.lengthInMeters}
  \uses{def:physics:phyx-mini-0738:phyxminiproblems-problemphyxmini0738-lengthquantity}
  Read a nonnegative physical length in metres
\end{definition}
\begin{proof}
  This is the declared model definition of length In Meters.
\end{proof}
\begin{definition}[signed Length In Meters]
  \label{def:physics:phyx-mini-0738:phyxminiproblems-problemphyxmini0738-signedlengthinmeters}
  \lean{PhyXMiniProblems.ProblemPhyXMini0738.signedLengthInMeters}
  \uses{def:physics:phyx-mini-0738:phyxminiproblems-problemphyxmini0738-signedlengthquantity}
  Read a signed Cartesian length component in metres
\end{definition}
\begin{proof}
  This is the declared model definition of signed Length In Meters.
\end{proof}
\begin{definition}[duration In Seconds]
  \label{def:physics:phyx-mini-0738:phyxminiproblems-problemphyxmini0738-durationinseconds}
  \lean{PhyXMiniProblems.ProblemPhyXMini0738.durationInSeconds}
  \uses{def:physics:phyx-mini-0738:phyxminiproblems-problemphyxmini0738-durationquantity}
  Read a physical duration in seconds
\end{definition}
\begin{proof}
  This is the declared model definition of duration In Seconds.
\end{proof}
\begin{definition}[speed Readout]
  \label{def:physics:phyx-mini-0738:phyxminiproblems-problemphyxmini0738-speedreadout}
  \lean{PhyXMiniProblems.ProblemPhyXMini0738.speedReadout}
  Read a speed in the selected length unit per selected time unit
\end{definition}
\begin{proof}
  This is the declared model definition of speed Readout.
\end{proof}
\begin{definition}[speed In Meters Per Second]
  \label{def:physics:phyx-mini-0738:phyxminiproblems-problemphyxmini0738-speedinmeterspersecond}
  \lean{PhyXMiniProblems.ProblemPhyXMini0738.speedInMetersPerSecond}
  \uses{def:physics:phyx-mini-0738:phyxminiproblems-problemphyxmini0738-speedreadout}
  Read an unsigned physical speed in metres per second
\end{definition}
\begin{proof}
  This is the declared model definition of speed In Meters Per Second.
\end{proof}
\begin{definition}[speed In Kilometers Per Hour]
  \label{def:physics:phyx-mini-0738:phyxminiproblems-problemphyxmini0738-speedinkilometersperhour}
  \lean{PhyXMiniProblems.ProblemPhyXMini0738.speedInKilometersPerHour}
  \uses{def:physics:phyx-mini-0738:phyxminiproblems-problemphyxmini0738-speedreadout}
  Read an unsigned physical speed in kilometres per hour
\end{definition}
\begin{proof}
  This is the declared model definition of speed In Kilometers Per Hour.
\end{proof}
\begin{definition}[signed Velocity In Meters Per Second]
  \label{def:physics:phyx-mini-0738:phyxminiproblems-problemphyxmini0738-signedvelocityinmeterspersecond}
  \lean{PhyXMiniProblems.ProblemPhyXMini0738.signedVelocityInMetersPerSecond}
  \uses{def:physics:phyx-mini-0738:phyxminiproblems-problemphyxmini0738-signedvelocityquantity}
  Read a signed velocity component in metres per second
\end{definition}
\begin{proof}
  This is the declared model definition of signed Velocity In Meters Per Second.
\end{proof}
\begin{definition}[acceleration In Meters Per Second Squared]
  \label{def:physics:phyx-mini-0738:phyxminiproblems-problemphyxmini0738-accelerationinmeterspersecondsquared}
  \lean{PhyXMiniProblems.ProblemPhyXMini0738.accelerationInMetersPerSecondSquared}
  \uses{def:physics:phyx-mini-0738:phyxminiproblems-problemphyxmini0738-accelerationquantity}
  Read an acceleration magnitude in metres per second squared
\end{definition}
\begin{proof}
  This is the declared model definition of acceleration In Meters Per Second Squared.
\end{proof}
\begin{definition}[degrees]
  \label{def:physics:phyx-mini-0738:phyxminiproblems-problemphyxmini0738-degrees}
  \lean{PhyXMiniProblems.ProblemPhyXMini0738.degrees}
  Convert a numerical degree readout into Mathlib's physical angle type
\end{definition}
\begin{proof}
  This is the declared model definition of degrees.
\end{proof}
\begin{definition}[Projectile Model]
  \label{def:physics:phyx-mini-0738:phyxminiproblems-problemphyxmini0738-projectilemodel}
  \lean{PhyXMiniProblems.ProblemPhyXMini0738.ProjectileModel}
  The ideal point-particle model used after the decoy is released
\end{definition}
\begin{proof}
  This is the declared model definition of Projectile Model.
\end{proof}
\begin{definition}[Figure Object]
  \label{def:physics:phyx-mini-0738:phyxminiproblems-problemphyxmini0738-figureobject}
  \lean{PhyXMiniProblems.ProblemPhyXMini0738.FigureObject}
  The objects visibly represented in image 738.png
\end{definition}
\begin{proof}
  This is the declared model definition of Figure Object.
\end{proof}
\begin{definition}[Figure Label]
  \label{def:physics:phyx-mini-0738:phyxminiproblems-problemphyxmini0738-figurelabel}
  \lean{PhyXMiniProblems.ProblemPhyXMini0738.FigureLabel}
  The two mathematical labels printed in the primary image
\end{definition}
\begin{proof}
  This is the declared model definition of Figure Label.
\end{proof}
\begin{definition}[Figure Point]
  \label{def:physics:phyx-mini-0738:phyxminiproblems-problemphyxmini0738-figurepoint}
  \lean{PhyXMiniProblems.ProblemPhyXMini0738.FigurePoint}
  The three points needed to interpret the displayed height and distance
\end{definition}
\begin{proof}
  This is the declared model definition of Figure Point.
\end{proof}
\begin{definition}[Planar Position]
  \label{def:physics:phyx-mini-0738:phyxminiproblems-problemphyxmini0738-planarposition}
  \lean{PhyXMiniProblems.ProblemPhyXMini0738.PlanarPosition}
  \uses{def:physics:phyx-mini-0738:phyxminiproblems-problemphyxmini0738-signedlengthquantity}
  A physical position in the vertical plane, with upward-positive height
\end{definition}
\begin{proof}
  This is the declared model definition of Planar Position.
\end{proof}
\begin{definition}[Planar Velocity]
  \label{def:physics:phyx-mini-0738:phyxminiproblems-problemphyxmini0738-planarvelocity}
  \lean{PhyXMiniProblems.ProblemPhyXMini0738.PlanarVelocity}
  \uses{def:physics:phyx-mini-0738:phyxminiproblems-problemphyxmini0738-signedvelocityquantity}
  A signed planar velocity, with upward-positive vertical component
\end{definition}
\begin{proof}
  This is the declared model definition of Planar Velocity.
\end{proof}
\begin{definition}[Airplane Decoy Figure]
  \label{def:physics:phyx-mini-0738:phyxminiproblems-problemphyxmini0738-airplanedecoyfigure}
  \lean{PhyXMiniProblems.ProblemPhyXMini0738.AirplaneDecoyFigure}
  \uses{def:physics:phyx-mini-0738:phyxminiproblems-problemphyxmini0738-lengthquantity, def:physics:phyx-mini-0738:phyxminiproblems-problemphyxmini0738-figureobject, def:physics:phyx-mini-0738:phyxminiproblems-problemphyxmini0738-figurelabel, def:physics:phyx-mini-0738:phyxminiproblems-problemphyxmini0738-figurepoint, def:physics:phyx-mini-0738:phyxminiproblems-problemphyxmini0738-planarposition}
  The labelled physical geometry and qualitative content of the supplied raster. The image is schematic, but coordinateOf interprets its three distinguished points in the physical horizontal/upward coordinate system
\end{definition}
\begin{proof}
  This is the declared model definition of Airplane Decoy Figure.
\end{proof}
\begin{definition}[Radar Decoy Release Setup]
  \label{def:physics:phyx-mini-0738:phyxminiproblems-problemphyxmini0738-radardecoyreleasesetup}
  \lean{PhyXMiniProblems.ProblemPhyXMini0738.RadarDecoyReleaseSetup}
  \uses{def:physics:phyx-mini-0738:phyxminiproblems-problemphyxmini0738-lengthquantity, def:physics:phyx-mini-0738:phyxminiproblems-problemphyxmini0738-durationquantity, def:physics:phyx-mini-0738:phyxminiproblems-problemphyxmini0738-accelerationquantity, def:physics:phyx-mini-0738:phyxminiproblems-problemphyxmini0738-projectilemodel, def:physics:phyx-mini-0738:phyxminiproblems-problemphyxmini0738-planarposition, def:physics:phyx-mini-0738:phyxminiproblems-problemphyxmini0738-planarvelocity, def:physics:phyx-mini-0738:phyxminiproblems-problemphyxmini0738-airplanedecoyfigure}
  Independent physical quantities for the passive release and subsequent flight. Neither the flight duration nor the release height is defined from an answer choice; governing laws below relate them to the observed impact
\end{definition}
\begin{proof}
  This is the declared model definition of Radar Decoy Release Setup.
\end{proof}
\begin{definition}[Matches Primary Airplane Decoy Figure]
  \label{def:physics:phyx-mini-0738:phyxminiproblems-problemphyxmini0738-matchesprimaryairplanedecoyfigure}
  \lean{PhyXMiniProblems.ProblemPhyXMini0738.MatchesPrimaryAirplaneDecoyFigure}
  \uses{def:physics:phyx-mini-0738:phyxminiproblems-problemphyxmini0738-lengthinmeters, def:physics:phyx-mini-0738:phyxminiproblems-problemphyxmini0738-signedlengthinmeters, def:physics:phyx-mini-0738:phyxminiproblems-problemphyxmini0738-radardecoyreleasesetup}
  Primary-image evidence. It records the objects, paths, arrows, labels, and spatial relations shown in 738.png; the release height is only identified as a vertical separation and is not assigned a numerical value
\end{definition}
\begin{proof}
  This is the declared model definition of Matches Primary Airplane Decoy Figure.
\end{proof}
\begin{definition}[Matches Problem Description]
  \label{def:physics:phyx-mini-0738:phyxminiproblems-problemphyxmini0738-matchesproblemdescription}
  \lean{PhyXMiniProblems.ProblemPhyXMini0738.MatchesProblemDescription}
  \uses{def:physics:phyx-mini-0738:phyxminiproblems-problemphyxmini0738-lengthinmeters, def:physics:phyx-mini-0738:phyxminiproblems-problemphyxmini0738-speedinkilometersperhour, def:physics:phyx-mini-0738:phyxminiproblems-problemphyxmini0738-degrees, def:physics:phyx-mini-0738:phyxminiproblems-problemphyxmini0738-radardecoyreleasesetup}
  Numerical readouts and qualitative idealization stated or conventionally implicit in the textbook projectile problem. No release-height value occurs in this structure
\end{definition}
\begin{proof}
  This is the declared model definition of Matches Problem Description.
\end{proof}
\begin{definition}[Uses Standard Near Earth Gravity]
  \label{def:physics:phyx-mini-0738:phyxminiproblems-problemphyxmini0738-usesstandardnearearthgravity}
  \lean{PhyXMiniProblems.ProblemPhyXMini0738.UsesStandardNearEarthGravity}
  \uses{def:physics:phyx-mini-0738:phyxminiproblems-problemphyxmini0738-accelerationinmeterspersecondsquared, def:physics:phyx-mini-0738:phyxminiproblems-problemphyxmini0738-radardecoyreleasesetup}
  The standard near-Earth gravitational calibration used in the calculation. This is a governing environmental datum, not a release-height assumption
\end{definition}
\begin{proof}
  This is the declared model definition of Uses Standard Near Earth Gravity.
\end{proof}
\begin{definition}[Has Physical Radar Decoy Parameters]
  \label{def:physics:phyx-mini-0738:phyxminiproblems-problemphyxmini0738-hasphysicalradardecoyparameters}
  \lean{PhyXMiniProblems.ProblemPhyXMini0738.HasPhysicalRadarDecoyParameters}
  \uses{def:physics:phyx-mini-0738:phyxminiproblems-problemphyxmini0738-lengthinmeters, def:physics:phyx-mini-0738:phyxminiproblems-problemphyxmini0738-durationinseconds, def:physics:phyx-mini-0738:phyxminiproblems-problemphyxmini0738-speedinmeterspersecond, def:physics:phyx-mini-0738:phyxminiproblems-problemphyxmini0738-accelerationinmeterspersecondsquared, def:physics:phyx-mini-0738:phyxminiproblems-problemphyxmini0738-radardecoyreleasesetup}
  Positivity and the rightward/downward branch depicted in the figure
\end{definition}
\begin{proof}
  This is the declared model definition of Has Physical Radar Decoy Parameters.
\end{proof}
\begin{definition}[Satisfies Passive Release Projectile Laws]
  \label{def:physics:phyx-mini-0738:phyxminiproblems-problemphyxmini0738-satisfiespassivereleaseprojectilelaws}
  \lean{PhyXMiniProblems.ProblemPhyXMini0738.SatisfiesPassiveReleaseProjectileLaws}
  \uses{def:physics:phyx-mini-0738:phyxminiproblems-problemphyxmini0738-durationquantity, def:physics:phyx-mini-0738:phyxminiproblems-problemphyxmini0738-signedlengthinmeters, def:physics:phyx-mini-0738:phyxminiproblems-problemphyxmini0738-durationinseconds, def:physics:phyx-mini-0738:phyxminiproblems-problemphyxmini0738-speedinmeterspersecond, def:physics:phyx-mini-0738:phyxminiproblems-problemphyxmini0738-signedvelocityinmeterspersecond, def:physics:phyx-mini-0738:phyxminiproblems-problemphyxmini0738-accelerationinmeterspersecondsquared, def:physics:phyx-mini-0738:phyxminiproblems-problemphyxmini0738-radardecoyreleasesetup}
  The passive-release and no-drag constant-gravity laws in coherent SI components: the decoy initially inherits the airplane's rightward and downward velocity; horizontal velocity then remains constant; and vertical displacement includes uniform downward gravitational acceleration. These are general trajectory relations for every physical elapsed duration. They contain neither the eliminated height formula nor an answer value
\end{definition}
\begin{proof}
  This is the declared model definition of Satisfies Passive Release Projectile Laws.
\end{proof}
\begin{definition}[computed Flight Time In Seconds]
  \label{def:physics:phyx-mini-0738:phyxminiproblems-problemphyxmini0738-computedflighttimeinseconds}
  \lean{PhyXMiniProblems.ProblemPhyXMini0738.computedFlightTimeInSeconds}
  \uses{def:physics:phyx-mini-0738:phyxminiproblems-problemphyxmini0738-degrees}
  The exact SI flight time obtained from the stated horizontal distance and the horizontal component of the stated airplane speed
\end{definition}
\begin{proof}
  This is the declared model definition of computed Flight Time In Seconds.
\end{proof}
\begin{definition}[computed Release Height In Meters]
  \label{def:physics:phyx-mini-0738:phyxminiproblems-problemphyxmini0738-computedreleaseheightinmeters}
  \lean{PhyXMiniProblems.ProblemPhyXMini0738.computedReleaseHeightInMeters}
  \uses{def:physics:phyx-mini-0738:phyxminiproblems-problemphyxmini0738-degrees, def:physics:phyx-mini-0738:phyxminiproblems-problemphyxmini0738-computedflighttimeinseconds}
  The exact unrounded release height predicted by the stated data and standard gravity after eliminating the flight time. The setup's independent physical height field is not definitionally equal to this scalar expression
\end{definition}
\begin{proof}
  This is the declared model definition of computed Release Height In Meters.
\end{proof}
\begin{definition}[Answer Choice]
  \label{def:physics:phyx-mini-0738:phyxminiproblems-problemphyxmini0738-answerchoice}
  \lean{PhyXMiniProblems.ProblemPhyXMini0738.AnswerChoice}
  Labels of the four displayed multiple-choice answers
\end{definition}
\begin{proof}
  This is the declared model definition of Answer Choice.
\end{proof}
\begin{definition}[height In Meters]
  \label{def:physics:phyx-mini-0738:phyxminiproblems-problemphyxmini0738-answerchoice-heightinmeters}
  \lean{PhyXMiniProblems.ProblemPhyXMini0738.AnswerChoice.heightInMeters}
  \uses{def:physics:phyx-mini-0738:phyxminiproblems-problemphyxmini0738-answerchoice}
  Height in metres printed beside each answer label
\end{definition}
\begin{proof}
  This is the declared model definition of height In Meters.
\end{proof}
\begin{definition}[recorded Dataset Answer]
  \label{def:physics:phyx-mini-0738:phyxminiproblems-problemphyxmini0738-recordeddatasetanswer}
  \lean{PhyXMiniProblems.ProblemPhyXMini0738.recordedDatasetAnswer}
  \uses{def:physics:phyx-mini-0738:phyxminiproblems-problemphyxmini0738-answerchoice}
  The answer label recorded in the supplied dataset metadata
\end{definition}
\begin{proof}
  This is the declared model definition of recorded Dataset Answer.
\end{proof}
\begin{definition}[Matches Displayed Whole Meter Height]
  \label{def:physics:phyx-mini-0738:phyxminiproblems-problemphyxmini0738-matchesdisplayedwholemeterheight}
  \lean{PhyXMiniProblems.ProblemPhyXMini0738.MatchesDisplayedWholeMeterHeight}
  \uses{def:physics:phyx-mini-0738:phyxminiproblems-problemphyxmini0738-lengthquantity, def:physics:phyx-mini-0738:phyxminiproblems-problemphyxmini0738-lengthinmeters, def:physics:phyx-mini-0738:phyxminiproblems-problemphyxmini0738-answerchoice}
  Agreement with a whole-metre display to nearest-metre precision
\end{definition}
\begin{proof}
  This is the declared model definition of Matches Displayed Whole Meter Height.
\end{proof}
\begin{definition}[Is Unique Closest Displayed Height]
  \label{def:physics:phyx-mini-0738:phyxminiproblems-problemphyxmini0738-isuniqueclosestdisplayedheight}
  \lean{PhyXMiniProblems.ProblemPhyXMini0738.IsUniqueClosestDisplayedHeight}
  \uses{def:physics:phyx-mini-0738:phyxminiproblems-problemphyxmini0738-lengthquantity, def:physics:phyx-mini-0738:phyxminiproblems-problemphyxmini0738-lengthinmeters, def:physics:phyx-mini-0738:phyxminiproblems-problemphyxmini0738-answerchoice}
  A displayed choice is strictly closer than every alternative
\end{definition}
\begin{proof}
  This is the declared model definition of Is Unique Closest Displayed Height.
\end{proof}
\begin{lemma}[release Height eq projectile Expression]
  \label{lem:physics:phyx-mini-0738:phyxminiproblems-problemphyxmini0738-releaseheight-eq-projectileexpression}
  \lean{PhyXMiniProblems.ProblemPhyXMini0738.releaseHeight_eq_projectileExpression}
  \uses{def:physics:phyx-mini-0738:phyxminiproblems-problemphyxmini0738-lengthinmeters, def:physics:phyx-mini-0738:phyxminiproblems-problemphyxmini0738-speedinmeterspersecond, def:physics:phyx-mini-0738:phyxminiproblems-problemphyxmini0738-accelerationinmeterspersecondsquared, def:physics:phyx-mini-0738:phyxminiproblems-problemphyxmini0738-radardecoyreleasesetup, def:physics:phyx-mini-0738:phyxminiproblems-problemphyxmini0738-matchesprimaryairplanedecoyfigure, def:physics:phyx-mini-0738:phyxminiproblems-problemphyxmini0738-hasphysicalradardecoyparameters, def:physics:phyx-mini-0738:phyxminiproblems-problemphyxmini0738-satisfiespassivereleaseprojectilelaws}
  Eliminating the positive flight duration from the two component equations gives the standard downward-launch formula h = d sin(theta) / cos(theta) + g d\textasciicircum{}2 / (2 (v cos(theta))\textasciicircum{}2). This conclusion is derived from the general laws and is not a premise field
\end{lemma}
\begin{proof}
  The conclusion follows from the governing relations and figure readouts named in the statement of release Height eq projectile Expression.
\end{proof}
% --- Archon declaration topology end ---
% --- Archon physics formalization source end ---
